\documentclass[../book.tex]{subfiles}

\begin{document}

    \section{一元函数积分学的应用 - 积分等式与积分不等式}

    \subsection{积分等式}

    \begin{enumerate}
        \item 用中值定理
        \item 用夹逼准则
        \item 用积分法: 恒等变形、换元法、分部积分法
    \end{enumerate}

    \subsection{积分不等式}

    \begin{enumerate}
        \item 用函数的单调性：首先将某一积分限（通常取上限）变量化，然后移项构造辅助函数，由辅助函数的单调性来证明不等式，此方法多用于所给条件为\textbf{\blueuline{$f(x)$在$[a, b]$上连续}}的情形
        \item 用拉格朗日中值定理：此方法多用于所给条件为\textbf{\blueuline{$f(x)$一阶可导}}且某一端点值比较简单（甚至为$0$）的题目
        \begin{itemize}
            \item 若区间$[a, b]$没有$x$，可将区间拆分为$[a, x]$和$[x, b]$，在两个区间上分别应用拉格朗日中值定理
        \end{itemize}
        \item 用泰勒公式：此方法多用于所给条件为\textbf{\blueuline{$f(x)$二阶可导}}且题中有简单函数值（甚至为$0$）的题目。用泰勒公式将被积函数$f(x)$展开成多项式再做
        \item 用积分法：被积函数是两项相乘的形式，可采用分部积分法
        \item 用牛顿-莱布尼茨公式
    \end{enumerate}

    \subsection{基础概念}

    \begin{enumerate}
    \end{enumerate}

    \subsection{结论}

    \begin{enumerate}
        \item $\lim\limits_{t\to 0^+} t^a\ln t = 0(a > 0)$
        \item 设$f(x)$在$[0, 1]$上连续，则$\lim\limits_{n\to \infty} \displaystyle\int_0^1 x^{n}f(x)\,\mathrm{d}x = 0$
    \end{enumerate}

    \subsection{定理}

    \begin{enumerate}
        \item 推广的积分中值定理: 设$f(x), g(x)$在$[a, b]$上连续，且$g(x)$在$[a, b]$上不变化，则存在$\xi \in (a, b)$，使得
        \[
            \int_a^b f(x)g(x)\,\mathrm{d}x = f(\xi)\int_a^b g(x)\,\mathrm{d}x
        \]
    \end{enumerate}

    \subsection{运算}

    \begin{enumerate}
        \item $\lim\limits_{n\to\infty} \displaystyle\int_0^1 \frac{x^{n+1}}{1+x}$
        \begin{analysisbox}[夹逼准则（放缩法）]
            由于在 $x\in[0,1]$ 上，
            \[
                0 \le \frac{x^{n+1}}{1+x} \le x^{n+1},
            \]
            因此有
            \[
                0 \le \int_0^1 \frac{x^{n+1}}{1+x}\,\mathrm{d}x
                \le \int_0^1 x^{n+1}\,\mathrm{d}x
                = \frac{1}{n+2}.
            \]

            当 $n\to\infty$ 时，
            \[
                \frac{1}{n+2} \to 0.
            \]

            由夹逼准则可得
            \[
                \lim_{n\to\infty}\int_0^1 \frac{x^{n+1}}{1+x}\,\mathrm{d}x
                = 0.
            \]
        \end{analysisbox}
    \end{enumerate}

    \subsection{公式}

    \begin{enumerate}

    \end{enumerate}

    \subsection{方法总结}

    \begin{enumerate}
        \item $\int_a^b f(x)\,\mathrm{d}x = F(b) - F(a)$
        \item $\int_a^b f(x)\,\mathrm{d}x = f(\xi)(b - a)$
        \item $\int_a^x f'(t)\,\mathrm{d}t = f(x) - f(a)$
        \item $f'(\xi)(x - a) = f(x) - f(a)$
        \item $f(x + 2) - f(x) = (\int_x^{x + 2} f'(t)\,\mathrm{d}t)'$
        \item $f(x) = f(0) + f'(0)x + \displaystyle\frac{f''(\xi)}{2}x^2$，若$f''(\xi) < 0$，则$f(x) \leq f(0) + f'(0)x$
        \item $|f(x)| = |f(0)| + |f'(0)x| + \left|\displaystyle\frac{f''(\xi)}{2}x^2\right|$
        \item 若$f(x) \leq g(x)$，则$\int_a^b f(x)\,\mathrm{d}x \leq \int_a^b g(x)\,\mathrm{d}x$
        \item $|\int u\,\mathrm{d}v| = |uv - \int v\,\mathrm{d}u| \leq |uv| - \int |v|\,\mathrm{d}u$
        \item $|\int_a^b| = |\int_a^c + \int_c^b| \leq |\int_a^c| + |\int_c^b|$
    \end{enumerate}

    \subsection{条件转换思路}

    \begin{enumerate}
        \item 设$f'(x)$在$[a, b]$上连续，且$f(a) = f(b) = 0$，证明：$|f(x)| \leq \displaystyle\frac{1}{2}\int_a^b |f'(x)|\,\mathrm{d}x$
        \begin{itemize}
            \item 逆运算: $|f(x)| \leq \displaystyle\frac{1}{2}\int_a^b |f'(x)|\,\mathrm{d}x \Rightarrow |f(x)| + |f(x)| \leq \displaystyle\int_a^b |f'(x)|\,\mathrm{d}x$
            \item 恒等变形: $|f(x)| = |f(x) - 0| = |f(x) - f(a)| = |f(x) - f(b)|$
        \end{itemize}
    \end{enumerate}

    \subsection{理解}

    \begin{enumerate}
        \item 设 $f(x)$ 在 $[1,2]$ 上连续，计算
        \[
            \lim_{n\to\infty}\int_{1}^{2} f(x)e^{-x^{n}}\,\mathrm{d}x.
        \]

        \begin{analysisbox}[Example]

            由于 $f(x)$ 在 $[1,2]$ 上连续，且 $e^{-x^n}$ 在 $[1,2]$ 上非负，
            由积分中值定理可知，存在 $\xi_n \in (1,2)$，使得
            \[
                \int_{1}^{2} f(x)e^{-x^n}\,\mathrm{d}x
                =
                f(\xi_n)\int_{1}^{2} e^{-x^n}\,\mathrm{d}x.
            \]

            因为 $f(x)$ 在 $[1,2]$ 上连续，
            故 $f(\xi_n)$ 有界。

            当 $x\in(1,2]$ 时，有
            \[
                e^{x^n} > x^n + 1 > 0,
            \]
            即
            \[
                \frac{1}{e^{x^n}} < \frac{1}{x^n + 1} < \frac{1}{x^n}
            \]

            于是得到
            \[
                0
                <
                \int_{1}^{2} e^{-x^n}\,\mathrm{d}x
                <
                \int_{1}^{2} x^{-n}\,\mathrm{d}x =
                \left.
                \frac{x^{1-n}}{1-n}
                \right|_{1}^{2}
                =
                \frac{1}{1-n}\left(2^{\,1-n}-1\right).
            \]

            又
            \[
                \lim_{n\to\infty} \frac{1}{1-n}\left(2^{\,1-n}-1\right) = 0
            \]

            由夹逼准则可得
            \[
                \lim_{n\to\infty}\int_{1}^{2} e^{-x^n}\,\mathrm{d}x = 0.
            \]

            又因 $f(\xi_n)$ 有界，
            故
            \[
                \lim_{n\to\infty}
                \int_{1}^{2} f(x)e^{-x^n}\,\mathrm{d}x
                =
                0.
            \]

        \end{analysisbox}
    \end{enumerate}

\end{document}
